%% ch02_background.tex — Chapter 2: Background and Related Work
%% HSP-Agile CCF-B Report, Ye Xujun, ECNU, July 2026

\chapter{Background and Related Work}
\label{ch:background}

Preventing specification-conformance defects requires tools that already exist
in formal methods and in LLM-based synthesis---but that have not been combined
for this purpose.
What follows is the background needed to see that gap clearly.
The exposition begins with SOFL/Agile-SOFL and the FSF first-match rule
(reusing \code{classify\_signal} from Chapter~\ref{ch:intro}), then reviews
LLM code generation, CEGIS, SMT solving, pattern detection, and mutation
testing.
A closing comparison table positions \textsc{SgDP} against prior systems.

% ─────────────────────────────────────────────────────────────────────────────
\section{SOFL and Agile-SOFL}
\label{sec:bg:sofl}
% ─────────────────────────────────────────────────────────────────────────────

\subsection{The SOFL Formal Method}

The prevention problem assumes a formal artefact that already names ordered
scenarios.
SOFL is that artefact's historical home.
Structured Object-Oriented Formal Language~\citep{liu1997sofl} is a
\emph{process-oriented} formal method synthesising Z notation, data flow
diagrams, and object-oriented structuring, bridging diagrammatic accessibility
with verification-grade semantics.
SOFL's development paradigm is \emph{stepwise refinement}
\citep{back1998refinement}: each refinement step must preserve abstract
specification properties, preventing specification drift.
Its \emph{process specification} enumerates named scenarios with guard
conditions and output bodies---the direct ancestor of FSF
(Section~\ref{sec:bg:fsf}) and the primary unit of the HSP-Agile pipeline.

\subsection{Agile-SOFL: Lightweight Formal Methods in Sprints}

\emph{Agile-SOFL} \citep{li2022agilesofl} reconciles SOFL with Agile
\citep{beck2001xp} by applying formal rigour to sprint-scoped behavioural
contracts: each sprint produces FSF specifications and code, with automatic
conformance checking, SMT test generation, and violation flagging before review.
The toolchain provides an FSF parser, a language server, and a conformance
checker; \citet{li2023iet} showed substantially higher fault detection than
manual review at guard boundaries.
HSP-Agile augments this with an LLM-driven synthesis loop automating
\emph{verification} and \emph{repair}, addressing the workflow-integration
barrier identified by \citet{woodcock2009formalsurvey}.

% ─────────────────────────────────────────────────────────────────────────────
\section{The Functional Specification Format (FSF)}
\label{sec:bg:fsf}
% ─────────────────────────────────────────────────────────────────────────────

\subsection{Syntax and Structure}

FSF is not a flat case statement.
Guards are ordered: the first match wins, and that order is part of the
contract.
The Functional Specification Format is the primary artefact of an Agile-SOFL
sprint \citep{li2022agilesofl, liu2014sofl}.
An FSF document describes a single function $t$ by its \emph{name},
\emph{parameter list}, \emph{external variables} (read/write), and an ordered
sequence of \emph{scenario clauses}:

\begin{center}
\small\ttfamily
\begin{minipage}{0.82\textwidth}
\obeylines
Function: \textit{name}(\textit{param$_1$}, \ldots, \textit{param$_n$})
~~ext rd \textit{x$_1$}, \textit{x$_2$}, \ldots
~~ext wr \textit{y$_1$}, \textit{y$_2$}, \ldots
~~Scenario S1: \textit{guard$_1$}
~~~~~~\textit{output definitions}
~~Scenario S2: \textit{guard$_2$}
~~~~~~\textit{output definitions}
~~Scenario Sk: others
~~~~~~\textit{default output definitions}
\end{minipage}
\end{center}

\noindent
Each scenario $S_i$ comprises (a) a Boolean \emph{guard condition} over
parameters and readable external variables, and (b) an \emph{assignment block}
specifying output values when the guard holds.
The final scenario may use the keyword \texttt{others} as an unconditional
catch-all and must appear last.

\subsection{Scenario Ordering Semantics}
\label{sec:bg:fsf:semantics}

The defining semantic feature of FSF is \emph{top-to-bottom priority}: the
first scenario whose guard is satisfied determines the function's output.
An input satisfying both $S_1$'s and $S_2$'s guards is handled exclusively
by $S_1$; $S_2$ applies only to inputs not already covered by higher-priority
scenarios.

The running example \code{classify\_signal} from Chapter~\ref{ch:intro}
makes the rule concrete (Listing~\ref{lst:fsf-spec}):

\begin{center}
\small\ttfamily
\begin{minipage}{0.85\textwidth}
\obeylines
Function: classify\_signal(level, threshold)
~~Scenario S1: level < 0
~~~~~~alert := "Error"
~~Scenario S2: level >= threshold
~~~~~~alert := "Critical"
~~Scenario S3: others
~~~~~~alert := "Normal"
\end{minipage}
\end{center}

\noindent
For input $(\mathit{level},\mathit{threshold}) = (-3,-10)$, both $S_1$'s
guard and $S_2$'s guard hold, yet FSF semantics require \texttt{"Error"}
because $S_1$ has higher priority.
An implementation that tests $\mathit{level} \ge \mathit{threshold}$ before
$\mathit{level} < 0$ incorrectly returns \texttt{"Critical"}---a
\emph{scenario-order violation} invisible to techniques that treat guards as
a flat Boolean formula.

Chapter~\ref{ch:formalization} returns to this partition in
Figure~\ref{fig:fsf-semantics}: $S_1$ claims the negative-level region even
when $S_2$'s threshold comparison would also hold; $S_3$ covers the remainder.
The precedence relation $S_1 \succ S_2 \succ S_3$ is structural, not derived
from the guards alone.

As a secondary pattern (not the running example), multi-output tasks such as
\code{classify\_risk(a,b,c)} exhibit the same first-match rule when a nested
guard ($a>10 \wedge b>5$) is a strict subset of a later guard ($a>5$).

\subsection{Formal Notation}

An FSF task is a triple $t = (\mathcal{S}_t,\, \mathcal{X}_t,\, \mathcal{Y}_t)$
where $\mathcal{X}_t$ and $\mathcal{Y}_t$ are input and output variable types,
and $\mathcal{S}_t = \langle s_1, s_2, \ldots, s_n \rangle$ is the ordered
scenario sequence ($s_1$ highest priority).
Each scenario $s_i = (g_i, \delta_i)$ comprises guard
$g_i : \mathcal{X}_t \to \{\mathtt{true}, \mathtt{false}\}$ and output
definition $\delta_i : \mathcal{X}_t \to \mathcal{Y}_t$.
The \emph{oracle} is:
\[
  g_t(x) \;=\; \delta_j(x)
  \quad\text{where}\quad
  j = \min\bigl\{i \mid g_i(x) = \mathtt{true}\bigr\}.
\]
A conforming implementation $f$ satisfies $f(x) = g_t(x)$ for all $x$;
any $x^*$ with $f(x^*) \neq g_t(x^*)$ is a \emph{counterexample}.

The oracle $g_t$ is \emph{not} equivalent to any unordered Boolean combination
of guards; a checker that ignores priority will accept ordering violations.
HSP-Agile's Z3-based checker encodes priority explicitly
(Chapter~\ref{ch:formalization}).
Unlike Dijkstra's guarded commands \citep{dijkstra1975guarded}, where
simultaneously true guards permit nondeterministic choice, FSF imposes
deterministic priority, yielding a unique executable oracle suitable for
automated test generation.

% ─────────────────────────────────────────────────────────────────────────────
\section{LLM-Assisted Code Generation}
\label{sec:bg:llm}
% ─────────────────────────────────────────────────────────────────────────────

\subsection{Foundation Models for Code}

Contest-style code generation optimises for high probability on hidden tests.
FSF synthesis needs something stricter: \emph{guaranteed} conformance to an
ordered scenario contract, not luck on a sampled suite
\citep{li2022alphacode}.
\citet{chen2021codex} introduced Codex and the \emph{pass@k} estimator
$\widehat{\mathrm{pass}@k} = 1 - \binom{n-c}{k}/\binom{n}{k}$ for $n$ samples
with $c$ fully passing (Codex: 28.8\% pass@1 on HumanEval).
\citet{openai2023gpt4} improved structured prompt-following, which helps with
scenario ordering and Python syntax---but does not by itself enforce
first-match precedence.

\subsection{Evaluation of Deployed Code Generation Tools}

\citet{dakhel2023copilot} evaluated GitHub Copilot systematically and found
that while Copilot handles straightforward problems, it struggles with precise
boundary handling, multi-condition logic, and implicit conventions---exactly
the logic FSF scenario guards encode.
Copilot suggestions were often \emph{semantically plausible} yet contained
subtle errors (off-by-one boundaries, incorrect operator precedence).
Semantic plausibility without formal correctness is the core challenge of
LLM-based formal specification implementation.

\subsection{Limitations of LLMs for Formal Specification Compliance}

Three systematic failure modes arise without additional guidance
\citep{dakhel2023copilot, lynn2024verifierloop}.
Each pattern is illustrated with a short fragment of the kind the checker
must catch on \code{classify\_signal}-style tasks:

\begin{enumerate}
  \item \textbf{Scenario-order blindness.}
        LLMs learn \texttt{if}--\texttt{elif} surface syntax but not
        specification-level priority ordering; overlapping guards bias the
        model toward ``natural'' code order rather than mandated precedence.
\begin{lstlisting}[language=Python,basicstyle=\ttfamily\footnotesize,frame=single]
# Wrong: threshold check before level < 0
if level >= threshold: return "Critical"
if level < 0: return "Error"   # never reached when both hold
\end{lstlisting}

  \item \textbf{Boundary hallucination.}
        Comparison operators (\texttt{>}, \texttt{>=}, \texttt{==}) are
        confused at boundaries; e.g.\ guard $\mathit{level} < 0$ implemented
        as $\mathit{level} \le 0$ is invisible to random testing but
        detectable by SMT boundary witnesses.
\begin{lstlisting}[language=Python,basicstyle=\ttfamily\footnotesize,frame=single]
# Wrong: off-by-one at the S1 boundary
if level <= 0: return "Error"   # accepts level == 0
\end{lstlisting}

  \item \textbf{Spurious completeness.}
        Hard-coded success values pass superficial tests but violate the
        specification on non-default inputs---corresponding to pattern RF07
        (\emph{unconditional success}, Section~\ref{sec:bg:patterns}).
\begin{lstlisting}[language=Python,basicstyle=\ttfamily\footnotesize,frame=single]
# Wrong: ignores the specification entirely
return "Normal"
\end{lstlisting}
\end{enumerate}

\noindent
Test-feedback repair loops \citep{chen2021codex} execute unit tests against
each candidate and include failures in subsequent prompts (the B2 baseline).
In the FSF context, tests must be \emph{derived from} the specification; fixed
test suites may miss scenario boundaries, allowing convergence on candidates
that pass all tests yet violate uncovered inputs.
HSP-Agile uses Z3 to \emph{synthesise} boundary-hitting witnesses guaranteed
to expose scenario-order violations.

% ─────────────────────────────────────────────────────────────────────────────
\section{Counterexample-Guided Synthesis (CEGIS)}
\label{sec:bg:cegis}
% ─────────────────────────────────────────────────────────────────────────────

\subsection{CEGIS and LLM-Augmented Variants}

A candidate can pass every unit test and still evaluate the wrong scenario
first.
Repair therefore needs counterexamples that name the violated region, not only
a failing assertion.
CEGIS decomposes synthesis into a \emph{synthesiser} and a \emph{verifier}
\citep{solarlezama2008sketching}: the verifier returns counterexamples that
constrain later iterations until success or budget exhaustion.
LLM-augmented variants replace template synthesis with free-form generation
while retaining formal verification
\citep{loughman2024dafnybench,lynn2024verifierloop}, but prior systems lack
FSF scenario-order semantics and pattern-guard pre-filters.

\subsection{HSP-Agile's CEGIS Adaptation}

HSP-Agile instantiates the \textsc{SgDP} CEGIS loop with three FSF-specialised
components.
The \emph{synthesiser} is an LLM prompted with the FSF specification and
accumulated repair context.
The \emph{verifier} is the Z3-based conformance checker executing the candidate
on SMT-generated test suite $\mathcal{D}(t)$.
The \emph{feedback aggregator} packages counterexamples
$(x^*, g_t(x^*), f_{c_k}(x^*))$ and high-severity requirement patterns into
structured repair context.

The departure from prior LLM+verifier systems is communicating \emph{both}
formal counterexamples \emph{and} pattern-guard violations, so logical errors
and structural anti-patterns can be corrected in the same iteration
(Figure~\ref{fig:cegis-loop}).

\begin{figure}[t]
  \centering
  \tikzinput{diagrams/tikz/cegis_loop}
  \caption{The HSP-Agile CEGIS loop.
           The LLM synthesiser produces candidate $c_k$ from FSF specification
           $\mathcal{S}_t$ and repair context.
           The conformance checker verifies $c_k$ against the Z3-generated test
           suite; failures become counterexamples.
           The feedback aggregator packages counterexamples and anti-patterns
           for the next iteration.
           If all tests pass and pattern guards are clear, $c_k$ is accepted.}
  \label{fig:cegis-loop}
\end{figure}

\noindent
Budget parameter $K$ limits repair iterations ($K = 3$ in the experiments);
if no conforming candidate is found, the system returns the highest-scoring
candidate as best-effort, consistent with budget-bounded CEGIS
\citep{barrett2018smt}.

% ─────────────────────────────────────────────────────────────────────────────
\section{SMT Solving and Formal Verification}
\label{sec:bg:smt}
% ─────────────────────────────────────────────────────────────────────────────

\subsection{Satisfiability Modulo Theories}

Satisfiability Modulo Theories (SMT) extends SAT with background theories
(linear arithmetic, arrays, uninterpreted functions, bitvectors)
\citep{barrett2018smt}.
An SMT solver returns a satisfying assignment or proves unsatisfiability;
the assignment is directly usable as a concrete test input.
\citet{demoura2008z3} introduced Z3, the de facto standard solver with
incremental solving capabilities exploited by HSP-Agile's conformance checker.

\subsection{Test Case Generation via SMT}

For FSF task $t = (\langle s_1, \ldots, s_n \rangle, \mathcal{X}_t, \mathcal{Y}_t)$,
the checker constructs test suite $\mathcal{D}(t)$ by solving, for each $s_i$:
\[
  \varphi_i \;=\; g_i(x) \;\wedge\; \bigwedge_{j < i} \neg g_j(x).
\]
Formula $\varphi_i$ characterises the \emph{effective} guard of $s_i$---inputs
in $s_i$'s region not covered by higher-priority scenarios.
A satisfying assignment $x_i^*$ must produce oracle output $\delta_i(x_i^*)$;
comparing candidate output $f(x_i^*)$ checks priority-order conformance.

Unlike QuickCheck \citep{claessen2000quickcheck}, which generates random inputs,
SMT generation is \emph{targeted}: it guarantees at least one witness per
strict scenario region, ensuring adjacent-scenario boundaries are always tested.

\subsection{Relationship to Full Formal Verification}

Dafny \citep{leino2010dafny} supports full deductive verification via
annotations (pre/post-conditions, loop invariants), proving correctness under
\emph{all} inputs---stronger than HSP-Agile's bounded checking.
However, annotation burden is prohibitive in Agile-SOFL sprints where engineers
have limited formal methods background and sprint velocity is paramount.
HSP-Agile occupies a middle ground: stronger than random testing via
SMT-targeted boundary witnesses, lighter than full verification with no
developer annotations.
Prevention guarantees are bounded by generated test-suite coverage
(Chapter~\ref{ch:formalization}).

% ─────────────────────────────────────────────────────────────────────────────
\section{Requirement-Level Error Pattern Detection}
\label{sec:bg:patterns}
% ─────────────────────────────────────────────────────────────────────────────

\subsection{Requirement-Level Patterns in FSF Implementations}

Static analysis tools such as FindBugs \citep{hovemeyer2004findbugs} encode
\emph{bug patterns}---recurring structural templates correlating with defects---
via visitor-based abstract interpretation \citep{cousot1977abstract}.
HSP-Agile's pattern guard extends this to \emph{requirements-level violations}:
structural signatures indicating the implementation cannot conform to the FSF
specification regardless of tested inputs.
Catalogue RF01--RF14 \citep{li2023iet} includes RF01 (missing precondition),
RF02 (wrong default branch), RF05 (swapped comparison), and RF07 (unconditional
success); see Appendix~\ref{app:patterns}.
High-severity patterns ($\mathcal{P}^H$) block acceptance and enter CEGIS repair
context even when all formal tests pass.

% ─────────────────────────────────────────────────────────────────────────────
\section{Mutation Testing for Specification Assessment}
\label{sec:bg:mutation}
% ─────────────────────────────────────────────────────────────────────────────

\subsection{Mutation Operators and Prevention Metrics}

Mutation testing \citep{demillo1978mutation,offutt1992mutation,ammann2016testing}
injects syntactic changes producing \emph{mutants}; survivors indicate inadequate
coverage.
HSP-Agile applies mutation at two levels.

\paragraph{Specification-level mutation} evaluates whether $\mathcal{D}(t)$
discriminates oracle $g_t$ variants:
\begin{itemize}
  \item \textbf{SNO} (Scenario Negate Once): negate one scenario's guard.
  \item \textbf{ORO} (Operator Replacement Ordinal): replace a relational operator
        ($>$ $\to$ $\geq$, $<$, $=$).
  \item \textbf{SPO} (Scenario Priority Order): swap adjacent scenario priorities.
\end{itemize}

\paragraph{Implementation-level mutation} evaluates candidate robustness:
\begin{itemize}
  \item \textbf{ICO} (Invert Condition Once): invert a Boolean sub-expression.
  \item \textbf{WRO} (Wrong Return Output): substitute a branch return value.
  \item \textbf{BCO} (Boundary Change Off-by-one): flip strict/non-strict comparison.
\end{itemize}

\noindent
HSP-Agile repurposes mutation as \emph{prevention assessment}: measuring how
well the pipeline \emph{prevents} non-conforming implementations from acceptance.

\begin{description}
  \item[Prevention Detection Rate (PDR).]
        Fraction of deliberately introduced faults (via implementation mutations)
        that the pipeline correctly rejects.

  \item[False Accept Rate (FAR).]
        Fraction of correct conforming candidates incorrectly rejected.
        Low FAR is essential for practical usability and developer trust.
\end{description}

\noindent
These metrics mirror sensitivity--specificity in binary classification
(Chapter~\ref{ch:formalization}).
\citet{okun2007mutation} used analogous detection/false-alarm rates for static
analysis evaluation, motivating PDR and FAR as primary prevention-quality measures.

% ─────────────────────────────────────────────────────────────────────────────
\section{Related Work}
\label{sec:bg:related}
% ─────────────────────────────────────────────────────────────────────────────

\subsection{Formal Methods in Software Engineering Practice}

\citet{woodcock2009formalsurvey} found that the most successful formal-method
deployments combine automated analysis with existing workflows rather than
replacing them wholesale.
They adopt formal methods in narrow, high-value layers.
Agile-SOFL follows this prescription; HSP-Agile automates synthesis and repair
without manual annotation effort.
Requirements-engineering research \citep{pohl2010requirements} emphasises
detecting ambiguities at the specification level.
HSP-Agile's pattern guard checks implementation structure against the FSF
specification before executing any test case.

\subsection{Approaches to LLM-Verified Code Synthesis}

Table~\ref{tab:related-work} compares closely related systems.

\paragraph{DafnyBench \citep{loughman2024dafnybench}.}
Evaluates LLM synthesis of Dafny programs with full deductive verification.
Repair relies on compiler error messages rather than semantic counterexamples;
no pattern guard pre-filter; evaluation uses pass@k without distinguishing
formal reasoning from statistical luck.

\paragraph{Lynn et al.\ TOSEM \citep{lynn2024verifierloop}.}
Verifier-in-the-loop repair closest to HSP-Agile, but without FSF scenario
ordering, pattern guards, or PDR/FAR prevention metrics; benchmarks lack
overlapping guard conditions where scenario-order violations arise.

\paragraph{Standard LLM (B1) and test-feedback (B2) baselines.}
B1 accepts the first LLM candidate without verification.
B2 adds iterative repair with fixed (non-SMT) unit tests---representative of
deployed coding assistants but without formal test generation or pattern guards.
The B2 vs.\ HSP-Agile comparison isolates SMT witness generation and pattern
guards over purely test-driven feedback.

\subsection{Structured Comparison}

\Needspace{14\baselineskip}
\begin{inlinetable}{Structured comparison of related approaches to LLM-assisted
           specification-conformant code generation.
           Checkmarks (\checkmark) indicate full support; dashes (--) absence;
           ``partial'' indicates limited support.
           F.\,Spec = formal specification; Prev.\,Metr.\ = prevention metrics.}
  \label{tab:related-work}
  \footnotesize
  \renewcommand{\arraystretch}{1.12}
  \begin{tabular*}{\linewidth}{@{\extracolsep{\fill}}l c c c c c@{}}
    \toprule
    \textbf{Approach}
      & \textbf{F.\,Spec}
      & \textbf{LLM}
      & \textbf{Repair}
      & \textbf{Guard}
      & \textbf{Prev.\,Metr.} \\
    \midrule
    DafnyBench \citep{loughman2024dafnybench}
      & \checkmark\ (Dafny)
      & \checkmark
      & partial
      & --
      & pass@k \\
    Lynn et al.\ TOSEM \citep{lynn2024verifierloop}
      & \checkmark
      & \checkmark
      & partial
      & --
      & correctness \\
    Standard LLM (B1)
      & --
      & \checkmark
      & --
      & --
      & pass rate \\
    Test-feedback (B2)
      & --
      & \checkmark
      & \checkmark\ (tests)
      & --
      & pass rate \\
    \midrule
    \textbf{HSP-Agile (ours)}
      & \checkmark\ \textbf{(FSF)}
      & \checkmark
      & \checkmark\ \textbf{(CEGIS)}
      & \checkmark
      & \textbf{PDR, FAR, Conf} \\
    \bottomrule
  \end{tabular*}%
\end{inlinetable}

\noindent
No prior approach combines formal FSF specification with scenario ordering,
LLM synthesis, CEGIS repair with SMT counterexamples, and a pattern guard layer.
Prior work evaluates via pass@k or aggregate correctness
\citep{chen2021codex, loughman2024dafnybench}, neither distinguishing false
acceptances from false rejections.
HSP-Agile introduces PDR and FAR to measure both error types directly.

% ─────────────────────────────────────────────────────────────────────────────
\section{Gap Summary}
\label{sec:bg:gap}
% ─────────────────────────────────────────────────────────────────────────────

LLM-assisted code generation has matured rapidly~\citep{chen2021codex,openai2023gpt4},
but 15--40\% of suggestions contain subtle semantic defects not caught by coarse
test suites~\citep{dakhel2023copilot}.
In formal specification settings these defects violate precisely-stated behavioural
contracts at boundary regions the specification characterises.
Prior SOFL studies showed early fault prevention reduces correction costs
\citep{li2022qrs,li2023iet}, but assumed human-written implementations; LLM
synthesis creates failure modes neither human-process prevention nor LLM
post-processing alone addresses.

Table~\ref{tab:related-positioning} positions \textsc{SgDP} on dimensions
critical for specification-conformance defect prevention.

\Needspace{3\baselineskip}
\paragraph{Agent-based code generation.}
Recent coding agent frameworks---OpenHands~\citep{wang2024opendevin},
SWE-agent~\citep{yang2024sweagent}, Codex CLI, and Claude Code---combine LLM
synthesis with tool use to resolve repository tasks.  They address a different
problem from \textsc{SgDP}: (i)~no specification oracle $g_t$;
(ii)~no prevention metrics (PDR/FAR); (iii)~orthogonal to the present
contribution---a stronger model could serve as Stage~1 synthesis backbone.
B5 (Reflexion-lite) captures the verbal-memory dimension most relevant to repair.
Full agent-oracle comparison is future work.

\begin{inlinetable}{Positioning \textsc{SgDP} against related approaches.  A tick indicates
  the property is satisfied; a dash indicates it is not a design goal.}
  \label{tab:related-positioning}
  \begin{tabular}{lccccc}
    \toprule
    & Spec-aware & Unified & Structured & Prevention & General \\
    Approach & witnesses & SpecIR & feedback IR & protocol & framework \\
    \midrule
    CEGIS (Solar-Lezama)          & ---  & --- & --- & ---  & \checkmark \\
    DafnyBench~\citep{dafnybench2024}     & \checkmark & --- & --- & --- & --- \\
    Self-Refine~\citep{madaan2023selfrefine}   & --- & --- & --- & --- & \checkmark \\
    Self-Debug~\citep{chen2023selfdebug}    & --- & --- & --- & --- & \checkmark \\
    Reflexion~\citep{shinn2023reflexion} (B5) & --- & --- & --- & --- & \checkmark \\
    CodeT~\citep{chen2022codet}           & --- & --- & --- & --- & --- \\
    OpenHands / SWE-agent       & --- & --- & --- & --- & \checkmark \\
    \textsc{SgDP} (this report)    & \checkmark & \checkmark & \checkmark & \checkmark & \checkmark \\
    \bottomrule
  \end{tabular}
\end{inlinetable}

Unlike VerifierLoop-style test-feedback systems~\citep{lynn2024verifierloop},
failing unit tests are not the primary signal: guard precedence is encoded as
SMT constraints~\citep{barrett2018smt}, and the LLM receives a typed failure
record that names the violated scenario.
Unlike DafnyBench-style proof-oriented pipelines~\citep{loughman2024dafnybench},
the model is not asked for annotations it systematically fails to produce.
Unlike Self-Refine / Self-Debug verbal loops, every repair step is grounded in
an executable specification oracle, and evaluation measures \emph{prevention}
(PDR/FAR), not only generation success.

The remaining gap that Table~\ref{tab:related-positioning} highlights is
therefore practical: a middle ground with specification-semantic awareness,
no proof obligation, and an evaluation protocol that asks whether
non-conforming code is kept out of release.
\Needspace{3\baselineskip}
That gap is what the formal problem statement in
Chapter~\ref{ch:formalization} makes precise.
